% ==========================================
% Archivo: sections/5_conclusions.tex
% ==========================================
\section{Conclusiones}

Del análisis de la dinámica multivariable del péndulo forzado se concluye que el sistema experimenta transiciones entre regímenes regulares y caóticos, fuertemente dependientes de la amplitud de forzamiento $F_D$. En los regímenes periódicos ($F_D \in \{0.50, 0.95\}$) y de re-estabilización ($F_D \in \{1.44, 4.00\}$), la dinámica está dominada por ciclos límite cerrados en el espacio de fases y un espectro de frecuencias discreto, garantizados por exponentes de Lyapunov negativos ($\lambda < 0$). Por el contrario, el régimen caótico ($F_D = 1.20$) desestabiliza las órbitas periódicas, generando atractores extraños de topología fractal evidenciados en las secciones de Poincaré de alta resolución, y provocando una divergencia exponencial de trayectorias vecinas ($\lambda > 0$). Asimismo, se demuestra que la ruptura de simetría en el amortiguamiento ($q^{(1)} \neq q^{(2)}$) arrastra al sistema hacia soluciones cuasiperiódicas ($\lambda \approx 0$), frustrando tanto el colapso asintótico como la divergencia caótica.

Cuantitativamente, la topología del diagrama de bifurcación confirma que la ruta hacia el caos sigue estrictamente una cascada de Feigenbaum por duplicación de periodo. El atractor fundamental de periodo 1 se desestabiliza en $F_D \approx 1.425$, bifurcándose en órbitas de periodo 2, y acelerando exponencialmente hacia órbitas de periodo 4 en $F_D \approx 1.46$. Esta acumulación de inestabilidades culmina en el punto crítico $F_D \approx  1.48$ (donde $\lambda = 0$), el cual acota la frontera de transición hacia un régimen aperiódico de bandas continuas para las variables cinemáticas $\theta$ y $\omega$.

Finalmente, la integración numérica del péndulo doble y del atractor de Lorenz corrobora la universalidad geométrica de la dinámica no lineal: la intrincada estructura de los atractores y la extrema sensibilidad a las condiciones iniciales no son artefactos del péndulo forzado, sino propiedades topológicas inherentes a los sistemas acoplados con más de dos grados de libertad en el espacio de fases.